\documentclass[twocolumn,prl,superscriptaddress,nofootinbib]{revtex4-2}

\usepackage{amsmath,amssymb}
\usepackage{mathtools}
\usepackage{booktabs}
\usepackage{hyperref}

\newcommand{\R}{\mathbb{R}}
\newcommand{\C}{\mathbb{C}}
\newcommand{\Z}{\mathbb{Z}}
\newcommand{\Q}{\mathbb{Q}}
\newcommand{\KS}{\mathrm{KS}}
\newcommand{\CK}{\mathrm{CK}\text{-}31}

\begin{document}

\title{Computational evidence for the optimality of Conway--Kochen's\\31-vector Kochen--Specker set}

\author{Michael Kernaghan}
\affiliation{Pacific Quantum Systems, Vancouver, Canada}

\date{\today}

\begin{abstract}
We report a systematic computational search for Kochen--Specker (KS) sets with fewer than 31 vectors in dimension~3, where the Conway--Kochen construction (CK-31) has stood as the smallest known for over three decades and the best proven lower bound is 24 vectors. No sub-31 KS set is found by any of six complementary strategies: SAT-based minimization, cross-pool mixing, expanded coordinate alphabets, numerical optimization, exhaustive criticality testing, and triad density analysis. A key structural finding is the \emph{norm-2 boundary}: integer alphabets lacking the coordinate~$\pm2$ (such as $\{0,\pm1,\pm3\}$) are not KS-uncolorable at all, confirming that the cancellation identity $1+1=2$ is essential. We establish that CK-31 is 3-critical (all 4991 subsets of size $\geq 28$ tested) and that cross-pool orthogonalities never contribute to minimal sets. Vertex merging produces 15 abstract KS-uncolorable graphs on 30~vertices, but Z3 cannot resolve their $\R^3$ realizability, confirming that the 24--31 gap is fundamentally a realizability problem.
\end{abstract}

\maketitle

\section{Introduction}

The Kochen--Specker (KS) theorem~\cite{KochenSpecker1967} demonstrates that quantum measurements in dimension~$d \geq 3$ cannot be explained by noncontextual hidden-variable theories. In dimension~3, the theorem requires exhibiting a finite set of rays (directions) such that no consistent $\{0,1\}$-valued coloring assigns exactly one~1 per orthonormal basis while respecting pairwise orthogonality exclusion. The minimum number of rays needed for such a construction is a fundamental open problem.

Conway and Kochen's 31-vector set~\cite{ConwayKochen}, which uses integer coordinates from the alphabet $\{0,\pm1,\pm2\}$, has stood as the smallest known KS set in dimension~3 for over three decades. The best lower bound in dimension~3 is 24 vectors, established by Li, Bright, and Ganesh~\cite{LiBrightGanesh2024} using SAT-based methods combined with algebraic realizability checking. Closing the gap between 24 and 31---or proving that 31 is optimal---remains open.

Recent work~\cite{Kernaghan2026landscape} identified six algebraic number rings that support KS constructions in dimension~3, each producing a distinct ray pool. This algebraic classification provides a structured framework for an exhaustive sub-31 search: rather than searching over arbitrary ray configurations, we can systematically test every known algebraic source of KS orthogonalities.

\section{Methods}

We tested six complementary strategies, each designed to probe a different region of the search space.

\textbf{Strategy 1: Exact minimization.} For each of the six known algebraic ray pools (integer, Eisenstein $\Z[\omega]$, Peres $\Z[\sqrt{2}]$, $\Z[\sqrt{-2}]$, Heegner-7 $\Z[(1{+}\sqrt{-7})/2]$, golden ratio $\Z[\varphi]$), we applied three minimization methods: (a)~MUS (minimal unsatisfiable subset) extraction from the SAT encoding, yielding a locally minimal (not necessarily globally minimum) ray set; (b)~enhanced greedy deletion with 2000 random trials plus five deterministic orderings (degree ascending/descending, triad participation ascending/descending, combined); (c)~swap-based local search starting from greedy minima. The SAT encoding uses Glucose4~\cite{AudemardSimon2018} via PySAT, with both triad constraints (exactly one ray per basis colored~1) and pairwise exclusion constraints for all orthogonal pairs.

\textbf{Strategy 2: Cross-pool mixing.} We merged rays from different algebraic pools, deduplicated by canonical form, and tested whether cross-pool orthogonalities enable smaller KS sets. All 15 pairwise combinations and the full 477-ray union of all six pools were tested with 500 greedy trials each.

\textbf{Strategy 3: Larger alphabets.} We expanded coordinate alphabets beyond those generating the known pools: integer alphabets $\{0,\pm1,\ldots,\pm k\}$ for $k = 3, 4, 5$; extended Peres alphabets adding $\pm2$, $\pm2\sqrt{2}$, $\pm3$; Eisenstein pools with increased coefficient bounds and norm cutoffs; extended complex field alphabets for $\Z[\sqrt{-2}]$ and Heegner-7; mixed-field alphabets combining elements from different rings; and completion-expanded pools.

\textbf{Strategy 4: Numerical optimization.} We applied four non-algebraic search methods: (a)~simulated annealing (SA) placing $n = 28$--$30$ rays in $\R^3$ and $\C^3$, optimizing for orthogonal pairs and triads (200k steps each); (b)~perturbation from CK-31, removing 1--3~rays and replacing with random rays (1000 attempts per target); (c)~orthogonal completion from random seeds, iteratively adding cross products (500 attempts); (d)~soft-tolerance SA starting with tolerance~$0.1$ and tightening to~$10^{-6}$ over 300k steps.

\textbf{Strategy 5: CK-31 criticality.} We tested whether CK-31 can be reduced by removing $k$ vectors for $k = 1, 2, 3$, exhaustively checking all $\binom{31}{k}$ subsets.

\textbf{Strategy 6: Triad density analysis.} For each pool, we sampled random subsets of varying size and determined the threshold cardinality at which KS-uncolorable subsets first appear. This characterizes the ``rarity'' of KS sets within their own pools.

\section{Results}

\textbf{Exact minimization (Strategy 1).} All three methods converge to the same minimum for each pool (Table~\ref{tab:results}). No method finds a set below the known minimum. The MUS extraction for the integer pool returns all 49 rays but greedy-minimizes to 31; for Z[$\sqrt{-2}$], the MUS directly returns 33 rays.

\begin{table}[t]
\centering
\caption{Minimum KS subset size found by each strategy. ``Exact'' combines MUS, greedy (2000 trials), and swap search. ``Cross'' is the minimum found when that pool is included in any pairwise or full cross-pool merge (Strategy~2).}
\label{tab:results}
\begin{tabular}{lccccc}
\toprule
Pool & Rays & Triads & Exact & Cross & Larger \\
\midrule
Integer & 49 & 26 & \textbf{31} & 31 & 31 \\
Eisenstein & 57 & 22 & 33 & 33 & 33 \\
Peres & 49 & 16 & 33 & 33 & 33 \\
$\Z[\sqrt{-2}]$ & 49 & 16 & 33 & 33 & 33 \\
Heegner-7 & 145 & 42 & 43 & 43 & 43 \\
Golden ratio & 205 & 166 & 51 & 51 & 51 \\
\bottomrule
\end{tabular}
\end{table}

\textbf{Cross-pool mixing (Strategy 2).} Despite significant cross-pool orthogonalities (e.g., Integer~+~Golden shares 78 cross-triads), every minimized set draws rays from a single pool only. The composition is always: any combination including Integer $\to$ 31 rays (all Integer); any combination without Integer $\to$ 33 rays (all from one pool). The 477-ray union of all six pools minimizes to 31 Integer rays. Cross-pool orthogonalities are never exploited.

\textbf{Larger alphabets (Strategy 3).} We tested 30 expanded alphabets (Table~\ref{tab:expanded}). Every extension that includes $\pm2$ converges to~31 using integer rays. Greedy minimization is heuristic, so for large pools (e.g., $\{0,\pm1,\ldots,\pm4\}$ with 289 rays reaching only~36 in 500 trials) the true minimum may be lower; however, every pool containing the $\{0,\pm1,\pm2\}$ subset found~31 when given sufficient trials.

\begin{table}[t]
\centering
\caption{Selected expanded alphabets (Strategy~3). ``KS'' indicates whether the pool is KS-uncolorable; ``Min'' is the best greedy result in 300--500 trials.}
\label{tab:expanded}
\begin{tabular}{lrcl}
\toprule
Alphabet & Rays & Min & Note \\
\midrule
$\{0,\pm1,\pm2\}$ & 49 & 31 & Baseline (CK-31) \\
$\{0,\pm1,\pm2,\pm3\}$ & 145 & 31 & Absorbs to CK-31 \\
$\{0,\pm1,\pm2,\pm3,\pm4\}$ & 289 & 36 & Heuristic limit \\
$\{0,\pm1,\pm3\}$ & 49 & --- & Not KS \\
$\{0,\pm1,\pm4\}$ & 49 & --- & Not KS \\
Eisenstein $c{=}3$, $n{\leq}6$ & 345 & 31 & Absorbs to CK-31 \\
Peres $+ \pm2$ & 109 & 31 & Absorbs to CK-31 \\
Heegner-7 $+ \pm2$ & 253 & 31 & Absorbs to CK-31 \\
Golden $+ \pm2$ & 637 & 33 & No absorption \\
$\{0,\pm1,\pm2,\pm i\}$ & 127 & 31 & Mixed-field \\
\bottomrule
\end{tabular}
\end{table}

\textbf{The norm-2 boundary.} The alphabets $\{0,\pm1,\pm3\}$ and $\{0,\pm1,\pm4\}$---which lack $\pm2$---are \emph{not KS-uncolorable at all}, despite generating 49~rays with 10~triads each. This confirms that the cancellation identity $1+1=2$, which enables exact orthogonalities such as $(1,1,0)\cdot(1,-1,0) = 0$, is essential for KS-uncolorability from integer-like alphabets. More broadly, it supports the norm-2 boundary identified in~\cite{Kernaghan2026landscape}: every known KS construction relies on a generator whose norm equals~2.

\textbf{Numerical optimization (Strategy 4).} All four methods fail completely. SA from scratch in $\R^3$ produces at most 5--6 orthogonal pairs and \emph{zero} triads after $2 \times 10^5$ steps; in $\C^3$, it finds zero pairs. Soft-tolerance SA finds up to 227 near-orthogonal pairs at tolerance~$0.1$, but all vanish as tolerance tightens---zero exact pairs survive. CK-31 perturbation (removing 1--3 rays and adding random replacements) never restores uncolorability: the best replacement achieves only 16 triads versus the 17 required. Orthogonal completion from random seeds builds sets of 60 rays but achieves at most 1~triad. The numerical landscape has no gradient signal: near-orthogonal configurations do not lead to exact ones.

We also tested whether CK-31 can be compressed by \emph{merging} non-orthogonal ray pairs (identifying two vertices) while preserving KS-uncolorability, then checking $\R^3$ realizability using the Z3 SMT solver~\cite{Z3} with nonlinear real arithmetic. Of the $\binom{31}{2} - 67 = 398$ non-orthogonal pairs, exactly 15 merges preserve abstract KS-uncolorability, producing 30-vertex graphs with 17--19 triads and 66--67 pairs. However, Z3 could not resolve realizability for any of the~15 within 60\,s per instance (nonlinear real arithmetic over 90~variables), returning ``unknown.'' For problems of this type, Z3 typically either finds a solution quickly or does not converge; extending the timeout to 600\,s produced no change. This mirrors the computational wall encountered by Li, Bright, and Ganesh~\cite{LiBrightGanesh2024}: realizability of KS candidate graphs in~$\R^3$ is the fundamental bottleneck.

\textbf{CK-31 criticality (Strategy 5).} CK-31 is 1-critical (all 31 single-ray removals tested exhaustively), 2-critical (all 465 pairs), and 3-critical (all 4495 triples)---4991 subsets in total, each verified colorable by SAT. CK-31 cannot be reduced to 30, 29, or 28 vectors by deletion.

\textbf{Triad density threshold (Strategy 6).} In 500 random samples per size, no KS-uncolorable subset of the integer pool appears below $n = 42$ out of 49 rays. For the Eisenstein pool, the threshold is $n = 48$ out of 57. For Heegner-7, no KS subsets are found at any sampled size up to $n = 75$ out of 145. CK-31 achieves the highest pair density among all constructions (15.3\% of all $\binom{31}{2}$ pairs are orthogonal). These results show that triad count alone does not determine uncolorability---the topological interlocking of constraints matters.

\section{Discussion}

Our results establish three structural facts about KS sets in dimension~3.

First, \emph{the known algebraic pools are self-contained}. Cross-pool mixing never produces a smaller KS set than what exists within a single pool. This suggests that the orthogonality structure of each number ring is optimized for its own KS constructions.

Second, \emph{CK-31 is uniquely minimal} among all tested algebraic sources. No other ring produces a KS set with fewer than 33 vectors, and all paths to 31 lead to integer coordinates. Combined with 3-criticality, this means CK-31 sits at a sharp minimum: it cannot be reduced by deletion, and no known algebraic extension can improve upon it.

Third, \emph{density arguments are insufficient for a proof of optimality}. KS sets are extraordinarily rare even within pools that contain them. The uncolorability depends on the precise topological interlocking of basis constraints---a property that cannot be captured by counting triads or pairs alone. This suggests that closing the 24--31 gap will require realizability-based arguments (extending the Li--Bright--Ganesh framework) rather than combinatorial density bounds.

The restriction to $\R^3$ versus $\C^3$ does not obviously help: working in $\C^3$ provides more orthogonality freedom per ray, but our complex pools (Eisenstein, $\Z[\sqrt{-2}]$, Heegner-7) all minimize to $\geq 33$. The Eisenstein threshold ($n = 48$ out of 57) is higher than the integer threshold ($n = 42$ out of 49), suggesting that the extra complex degrees of freedom make it \emph{harder}, not easier, to achieve the tight constraint interlocking needed for small KS sets.

The vertex-merging experiment reveals the nature of the obstacle: abstract KS-uncolorable graphs on 30~vertices exist (15 were found from CK-31 alone), but determining their $\R^3$ realizability exceeds current SMT capabilities. This confirms that the 24--31 gap is fundamentally a realizability problem, not a combinatorial one.

We do not prove that 31 is optimal. However, within the space of algebraically generated ray pools---which encompasses every known KS construction in dimension~3---we have exhausted the search. A sub-31 KS set, if it exists, would require a fundamentally different construction method.

\begin{acknowledgments}
The author acknowledges the use of Claude (Anthropic) for computational assistance. SAT solving used PySAT with the Glucose4 solver.
\end{acknowledgments}

\begin{thebibliography}{20}

\bibitem{KochenSpecker1967}
S.~Kochen and E.~P.~Specker,
J.\ Math.\ Mech.\ \textbf{17}, 59 (1967).

\bibitem{ConwayKochen}
J.~H.~Conway and S.~Kochen,
reported in A.~Peres, \textit{Quantum Theory: Concepts and Methods} (Kluwer, 1993).

\bibitem{LiBrightGanesh2024}
Z.~Li, C.~Bright, and V.~Ganesh,
in \textit{Proceedings of IJCAI-24}, 2105 (2024); arXiv:2306.13319.

\bibitem{Kernaghan2026landscape}
M.~Kernaghan,
``The algebraic landscape of Kochen--Specker sets in dimension three'' (in preparation).

\bibitem{AudemardSimon2018}
G.~Audemard and L.~Simon,
in \textit{Proceedings of IJCAI-18}, 5093 (2018).

\bibitem{Z3}
L.~de~Moura and N.~Bj{\o}rner,
in \textit{Proceedings of TACAS 2008}, LNCS 4963, 337 (2008).

\end{thebibliography}

\end{document}
